\section{What is a weak-field in General Relativity?}\label{sec:what-is-a-weak-field-in-general-relativity?}

The theory of general relativity is a field theory of gravitation that describes gravity as a result of the curvature of spacetime caused by massive objects.
In many astrophysical scenarios, the gravitational fields involved are relatively weak compared to the extreme conditions found near black holes or neutron stars.
In such cases, the weak-field approximation can be employed to simplify the complex equations of general relativity.

The weak-field approximation is particularly useful in scenarios such as gravitational wave propagation in the far-field region, where the gravitational waves are weak and can be treated as small perturbations on a flat background spacetime.
This approximation is also applicable in the study of gravitational lensing, orbital dynamics of binary systems, and the analysis of gravitational wave signals detected by observatories like LIGO and Virgo.
The latter being the main focus of this document.


\section{Weak-field metric}\label{sec:weak-field-metric}
The weak-field approximation assumes that there exists a global inertial frame where the metric tensor $g_{\mu\nu}$ can be expressed as a small perturbation $h_{\mu\nu}$ around the flat Minkowski metric $\eta_{\mu\nu}$:

\begin{equation}
    g_{\mu\nu} = \eta_{\mu\nu} + h_{\mu\nu}, \quad |h_{\mu\nu}| \ll 1\label{eq:weak-field-metric}
\end{equation}

This allows for a linearization of Einstein's field equations, making them more tractable for analytical and numerical solutions.

Einstein's theory of relativity, returns to Newtonian gravity in the limit of weak and static fields and the particles in such field are moving slowly compared to the speed of light.
Here we just restricted ourselves to the weak-field limit, and not the static case.
Therefore, the weak-field approximation is still a general relativistic treatment of gravity, but it simplifies the equations by assuming that the gravitational field is weak enough that the spacetime curvature is small.

Such smallness of fields usually happens when we are far away from the sources of gravity, this condition is usually called the far-field limit.
If (as it is usually the case) the far-field limit is to be a flat spacetime, then the metric can be approximated as in equation~\ref{eq:weak-field-metric}.
Which is called the asymptotically flat spacetime.

Formally, the weak-field approximation can be regarded as a field theory in which, $\eta_{\mu\nu}$ is a background metric; the gravitational field generated by the matter does not backreact on the source; $h_{\mu\nu}$ is the main field and transforms as a tensor on flat spacetime under Lorentz transformations.

\subsection{Change of Coordinates}\label{subsec:change-of-coordinates}

Consider a Lorentz transformation of coordinates ($\Lambda^{\text{T}}\eta\Lambda = \eta$):
\begin{equation}
    x^{\mu} = \Lambda^{\mu}_{\ \nu} x'^{\nu}\label{eq:coord-transformation-general-def}
\end{equation}

Then for the metric we have:
\begin{equation}
    g_{\mu'\nu'} = \frac{\partial x^{\alpha}}{\partial x'^{\mu'}} \frac{\partial x^{\beta}}{\partial x'^{\nu'}} g_{\alpha\beta} = \Lambda^{\alpha}_{\ \mu'} \Lambda^{\beta}_{\ \nu'} (\eta_{\alpha\beta} + h_{\alpha\beta}) = \eta_{\mu'\nu'} + \Lambda^{\alpha}_{\ \mu'} \Lambda^{\beta}_{\ \nu'} h_{\alpha\beta}\label{eq:metric-transformation-general-def}
\end{equation}

that implies that the following substitutions leaves the linearized theory invariant:
\begin{equation}
    \eta_{\mu\nu} \rightarrow \eta_{\mu'\nu'}, \quad h_{\mu\nu} \rightarrow h_{\mu'\nu'} = \Lambda^{\alpha}_{\ \mu'} \Lambda^{\beta}_{\ \nu'} h_{\alpha\beta}\label{eq:weak-field-metric-lorentz-invariance}
\end{equation}

\subsection{Infinitesimal Diffeomorphism Invariance}\label{subsec:infinitesimal-diffeomorphism-invariance}

Consider an infinitesimal coordinate transformation:
\begin{equation}
    x^{\mu} \rightarrow x^{\mu '} = x^{\mu} + \xi^{\mu}(x)\label{eq:infinitesimal-diffeomorphism}
\end{equation}

where,
\begin{equation}
    |\partial_{\beta}\xi^\alpha| \sim |h_{\mu\nu}| \ll 1.\label{eq:infinitesimal-diffeomorphism-condition}
\end{equation}

where one uses the Taylor expansion of a matrix $A = (1 + \delta \xi)$ for the inverse $(1 + \delta \xi)^{-1}\approx 1 - \delta \xi$.
To linear order in $h$ and in $\partial \xi$ the metric change is:
\begin{align}
    g_{\mu'\nu'} &= \frac{\partial x^{\mu}}{\partial x'^{\mu'}} \frac{\partial x^{\nu}}{\partial x'^{\nu'}} g_{\mu\nu} \nonumber  = (\delta^\mu_{\mu'} - \partial_{\mu'}\xi^\mu)(\delta^\nu_{\nu'} - \partial_{\nu'}\xi^\nu)(\eta_{\mu\nu} + h_{\mu\nu}) \nonumber \\
    & = \eta_{\mu'\nu'} + h_{\mu'\nu'} - 2\partial_{(\mu'}\xi_{\nu')}\label{eq:metric-transformation-infinitesimal-diffeomorphism}
\end{align}

Therefore, the metric is invariant under any infinitesimal coordinate transformation of the perturbed metric:

\begin{equation}
    \boxed{
        h_{\mu\nu} \rightarrow h_{\mu'\nu'} = h_{\mu\nu} +2\partial_{(\mu}\xi_{\nu)}
    }\label{eq:invariance-under-infinitesimal-diffeomorphism}
\end{equation}


\section{Derivations and Properties}\label{sec:derivations-and-properties}

\subsection{Iverse Metric}\label{subsec:iverse-metric}
To find the inverse metric $g^{\mu\nu}$, we start from the definition:
\begin{equation}
    g_{\mu\alpha} g^{\alpha\nu} = \delta_{\mu}^{\nu}\label{eq:inverse-metric-definition}
\end{equation}

Assuming a perturbative expansion for the inverse metric:
\begin{equation}
    g^{\mu\nu} = \eta^{\mu\nu} + k^{\mu\nu}\label{eq:inverse-metric-perturbation}
\end{equation}
we calculate the defined condition~\ref{eq:inverse-metric-definition},
\begin{align}
    g^{\alpha\mu}g_{\mu\beta} &= (\eta^{\alpha\mu} + k^{\alpha\mu})(\eta_{\mu\beta} + h_{\mu\beta}) \nonumber \\
    &= \delta^{\alpha}_{\beta} + \eta^{\alpha\mu} h_{\mu\beta} + k^{\alpha\mu}\eta_{\mu\beta} = \delta^{\alpha}_{\beta}\nonumber \ \ \ \cdot\eta^{\gamma\beta}  \\
    &\Rightarrow k^{\alpha\gamma} = - \eta^{\gamma\beta}\eta^{\alpha\mu} h_{\mu\beta} \eqqcolon -h^{\alpha\gamma}.\label{eq:inverse-metric-calculation-step}
\end{align}

\subsection{Other Quantities in Weak-field Approximation}\label{subsec:other-quantities-in-weak-field-approximation}

In the weak-field approximation, other important quantities can be expressed in terms of the perturbation $h_{\mu\nu}$.
We give a few final results are here, and the detailed derivations can be found in standard textbooks on general relativity.
\begin{equation}
    g^{\mu\nu} = \eta^{\mu\nu} + h^{\mu\nu}\mathcal{O}(h^2) \label{eq:inverse-metric-weak-field}
\end{equation}
And the christoffel symbols are:
\begin{equation}
    \Gamma^{\mu}_{\alpha\beta} = \frac12 \eta^{\mu\lambda} \left(
                                                               \partial_{\alpha} h_{\lambda\beta} + \partial_{\beta}h_{\lambda\alpha} - \partial_{\lambda}h_{\alpha\beta}
    \right) + \mathcal{O}(h^2)\label{eq:christoffel-symbols-weak-field}
\end{equation}
and for the Riemann tensor:
\begin{equation}
    R_{\mu\nu\alpha\beta} = \frac12 \left(
                                        \partial_{\alpha}\partial_{\nu} h_{\mu\beta} + \partial_{\beta}\partial_{\mu} h_{\nu\alpha} - \partial_{\alpha}\partial_{\mu}h_{\nu\beta} - \partial_{\beta}\partial_{\nu} h_{\mu\alpha}
    \right)\label{eq:riemman-tensor-weak-field}
\end{equation}
then for the Ricci tensor we have:
\begin{equation}
    R_{\mu\nu} = \partial^{\rho}\partial_{(\mu}h_{\nu)\rho} -\frac12 \partial^{\rho}\partial_{\rho}h_{\mu\nu} - \partial_{\mu}\partial_{\nu}h\label{eq:ricci-tensor-weak-field}
\end{equation}
where $h \coloneqq h^\mu_\mu$, and $\partial^\mu \coloneq \eta^{\mu\rho}\partial_\rho$.
And for the Einstein tensor:
\begin{equation}
    G_{\mu\nu} = \partial^{\rho}\partial_{(\mu} h_{\nu)\rho} -\frac12 \partial^{\rho}\partial_{\rho}h_{\mu\nu} - \partial_{\mu}\partial_{\nu}h - \frac12 \eta_{\mu\nu}\left(\partial^{\rho}\partial^{\sigma}h_{\rho\sigma} - \partial^{\rho}\partial_{\rho}h\right)\label{eq:einstein-tensor-weak-field}
\end{equation}

These equations may not look too informative, but we can already derive two important properties that will be important later on.
First, note that the second term on the right-hand side of equation~\ref{eq:ricci-tensor-weak-field} is just a wave operator (d'Alembertian) acting on the perturbation $h_{\mu\nu}$.
This indicates that gravitational perturbations propagate as waves in the weak-field limit.

Wave operators are nice in the sense that they lead to manifestly hyperbolic equations, which have well-posed initial value problems.


\section{Weak-field Solutions}\label{sec:weak-field-solutions}

